\documentclass[12pt]{article}
\usepackage[T1]{fontenc}
\usepackage[utf8]{inputenc}
\usepackage{lmodern}
\usepackage{textcomp}
\usepackage{enumitem}
\usepackage{geometry}
\geometry{margin=1in}
\begin{document}
\begin{center}
Answer Key - Physics - Undergraduate Level Assessment 7 \
\small (Answers are ordered exactly as in the question paper.)
\end{center}

\section*{Multiple Choice}
\begin{enumerate}[label=\arabic*.]
\item \textbf{Answer:} A
\\ \textit{Options:} A) gamma rays; B) microwaves; C) ultraviolet radiation; D) visible light
\\ \textit{Note:} Electromagnetic radiation emitted from a nucleus is most likely to be in the form of gamma rays because gamma rays are high-energy photons released during nuclear decay processes, reflecting the transition of the nucleus to a lower energy state.
\item \textbf{Answer:} D
\\ \textit{Options:} A) 0.25 W; B) 0.5 W; C) 1 W; D) 4 W
\\ \textit{Note:} Doubling the voltage across the resistor increases the power dissipation by a factor of four, as power is proportional to the square of the voltage (P = V²/R), leading to a new rate of energy dissipation of 4 W.
\item \textbf{Answer:} B
\\ \textit{Options:} A) 0.50 c; B) 0.60 c; C) 0.70 c; D) 0.80 c
\\ \textit{Note:} The correct answer is 0.60 c because, according to the principles of special relativity, length contraction occurs when an object is observed moving relative to an observer, and the contracted length (0.80 m) can be calculated using the formula L = L₀√(1 - v²/c²), where L₀ is the proper length (1.00 m).
\item \textbf{Answer:} A
\\ \textit{Options:} A) a dye laser; B) a helium-neon laser; C) an excimer laser; D) a ruby laser
\\ \textit{Note:} A dye laser is the correct answer because it can be tuned to emit a wide range of wavelengths across the visible spectrum, making it ideal for versatile spectroscopy applications.
\item \textbf{Answer:} C
\\ \textit{Options:} A) 1.2 x $10^9$ m/s; B) 3.0 x $10^8$ m/s; C) 1.5 x $10^8$ m/s; D) 1.0 x $10^8$ m/s
\\ \textit{Note:} The speed of light in a dielectric material is given by the equation c/n, where c is the speed of light in a vacuum (approximately 3 x $10^8$ m/s) and n is the refractive index, which is the square root of the dielectric constant; thus, with a dielectric constant of 4.0, the speed of light is 3 x $10^8$ m/s divided by 2 (the square root of 4), resulting in 1.5 x $10^8$ m/s.
\end{enumerate}

\section*{True / False}
\begin{enumerate}[label=\arabic*.]
\setcounter{enumi}{5}
\item \textbf{Answer:} False
\\ \textit{Options:} A) True; B) False
\\ \textit{Note:} The statement is false because the Sun's energy primarily arises from the fusion of four hydrogen nuclei into one helium nucleus, not just two hydrogen atoms.
\item \textbf{Answer:} False
\\ \textit{Options:} A) True; B) False
\\ \textit{Note:} The statement is false because electrons in metals are delocalized and can move freely, resulting in higher mean kinetic energy that can exceed kT due to their ability to participate in thermal conduction and collective excitations.
\item \textbf{Answer:} False
\\ \textit{Options:} A) True; B) False
\\ \textit{Note:} A helium-neon laser, while useful for some applications, is limited to specific wavelengths (primarily red) and lacks the tunability needed for versatile spectroscopy across the entire visible spectrum, making it less suitable compared to other laser types like tunable dye lasers or solid-state lasers.
\item \textbf{Answer:} False
\\ \textit{Options:} A) True; B) False
\\ \textit{Note:} The statement is false because de Broglie proposed that the linear momentum of a free massive particle is related to its wavelength through Planck's constant, not Boltzmann's constant.
\item \textbf{Answer:} False
\\ \textit{Options:} A) True; B) False
\\ \textit{Note:} In the diamond structure of elemental carbon, the nearest neighbors of each carbon atom are not located at the corners of a cube, but rather at the vertices of a tetrahedron formed by four other carbon atoms.
\end{enumerate}

\section*{Long Form Response}
\begin{enumerate}[label=\arabic*.]
\setcounter{enumi}{10}
\item \textbf{Answer:} The Pauli exclusion principle states that no two fermions, such as electrons, can occupy the same quantum state simultaneously. In the case of the helium atom, which has two electrons, this principle leads to the formation of different quantum states known as ortho- and para-helium, depending on the alignment of their spins. In the ortho-state, both electrons have parallel spins, which results in a higher total spin and, consequently, higher energy due to increased electron-electron repulsion. Conversely, the para-state features anti-parallel spins, leading to a lower total spin and reduced energy, as the electrons can occupy states that minimize their mutual repulsion. Thus, the energy difference between these two states can be directly attributed to the constraints imposed by the Pauli exclusion principle on electron configurations and their resulting interactions.
\\ \textit{Note:} correct because it accurately describes how the Pauli exclusion principle leads to different spin alignments in helium's electrons, resulting in varying energy levels for the ortho-state and para-state due to differences in electron-electron repulsion.
\item \textbf{Answer:} The emission spectra of doubly ionized lithium (Li++) and hydrogen exhibit notable similarities due to their analogous electronic structures, both containing a single electron orbiting a nucleus. In both cases, the wavelengths of emitted light correspond to transitions between quantized energy levels, which can be described by the Rydberg formula. However, for Li++, the wavelengths are scaled down by a factor of 9 compared to hydrogen, owing to the greater nuclear charge, which results in a stronger electrostatic attraction between the nucleus and the electron. This scaling has significant physical implications, indicating that as the atomic number increases, the energy levels become more closely spaced, leading to shorter wavelength emissions. Thus, while the fundamental nature of the spectra remains similar, the differences in wavelength highlight the effects of increased nuclear charge on electron behavior in multi-electron systems.
\\ \textit{Note:} The answer correctly identifies that both Li++ and hydrogen have a single electron and that their emission spectra can be described by the Rydberg formula, while also noting the significant effect of Li++'s higher nuclear charge, which scales the wavelengths of emitted light down by a factor of 9 compared to hydrogen.
\end{enumerate}

\vfill
\noindent\textit{End of Answer Key}
\end{document}